\chapter[Computational methods in BSSs]{Computational methods in \acp{bss}}
\chaptermark{Computational methods in BSSs}
\label{ch:computational_methods_in_bss}

The integration of computational methods into \acp{bss} marks a fundamental paradigm shift in urban mobility management. As these systems have progressed from early, small-scale bike fleets to today's interconnected large-scale systems, the sheer volume and complexity of generated data has surpassed the limits of traditional heuristic or manual decision-making approaches. Today, \acp{bss} operate as complex physical-digital systems in which thousands of vehicles, user interactions, and environmental conditions evolve continuously. In this context, computational methods have evolved from serving as supplementary analytical tools to becoming central to the operations of modern shared micromobility systems.

The application of computational methods in this domain is multifaceted, addressing challenges that range from the immediate (e.g., predicting demand for the next 15 minutes) to the long-term (e.g., designing infrastructure for the next decades). While the core contributions of this thesis focus specifically on station planning, and battery and maintenance dynamics modelling—topics which will be detailed in Chapters~\ref{ch:station_planning_in_barcelona}, \ref{ch:mobility_and_battery_modeling_bicing}, and \ref{ch:failure_forecasting_bicing}—it is essential to situate these contributions within the broader landscape of computational methods applications in \acp{bss}. This Chapter therefore provides an integrated overview of the data generated by \acp{bss} and surveys the principal fields in which computational methods are currently applied or actively researched.

\section[Data landscape]{Data landscape}

A clear understanding of the data underpinning computational methods in \acp{bss} is essential before delving into their applications. The core component of this landscape is trip-level mobility data, which typically include the \ac{od} stations of each ride, timestamps marking the start and end of the journey, and a unique bicycle identifier~\cite{FaghihImani2017, ZaltzAustwick2013}. In services that operate with user accounts, anonymized user identifiers may also be available, enabling longitudinal analyses of individual behavior while preserving privacy~\cite{FaghihImani2017}. Building on these basic records, some systems also capture detailed \ac{gps} trajectories, which provide the full sequence of points describing the traveled path~\cite{Talavera_Garcia2021}. When available, these traces expand the analytical possibilities by allowing the study of how individuals navigate the street network and how infrastructure influences route choice.

Beyond trip-level records, docked \acp{bss} also monitor station metadata and station state. Metadata typically include attributes such as station location and capacity, while station state tracks time-varying availability of bikes and docks~\cite{FaghihImani2014, CortezOrdonez2024}. Closely related with stations, many systems document rebalancing activities, noting the timing, locations, and quantities of bicycles relocated.

At the vehicle level, newer-generation fleets can generate additional operational data that extend beyond geolocation. Using smart locks and onboard sensors, bicycles can report lock/unlock bike events, and motion and attitude measurements (e.g., accelerometer and gyroscope) that enable monitoring and alerts such as fallen-bike detection~\cite{Xue2020}. Moreover, bicycles have been used as mobile sensing platforms for multiple purposes. For example, vibration signals have been leveraged for pavement-quality assessment~\cite{Qian2020Predicting}, while onboard video streams have been used for lane recognition~\cite{Kaundanya2024}. Furthermore, air-quality sensors have been mounted on bicycles to monitor pollutant exposure at fine spatial scales~\cite{Wu2020}.

Another source of information is maintenance operations data, which document repairs, component replacements, and other interventions performed on the fleet. These records are generally not publicly available, as they relate to internal operational processes and may contain sensitive information about system performance. Nevertheless, operators systematically collect them. Evidence of this is that planning guidelines explicitly recommend maintaining detailed maintenance logs to support effective system management~\cite{Yanocha2018}, and several studies have analysed such datasets when operators have made them accessible for research purposes~\cite{deBrito2022, Gregoire2018}. Maintenance data are highly valuable, as they can reveal recurrent failure patterns, inform maintenance strategies, and help optimize workshop operations and spare-parts procurement.

The increasing electrification of \ac{bss} fleets brings with it data on battery and sensor performance, such as state-of-charge, voltage, current, temperature, and charging events~\cite{Waldner2025}. These indicators equip operators to monitor and manage fleet energy levels, prioritize necessary interventions, and, at a research level, facilitate analyses of energy use patterns, battery health, and optimal charging strategies. 

Together, these datasets enable a detailed understanding of both organic user flows and operator-led adjustments, supporting evidence-based operational control.

Interestingly, a standardized subset of these system-generated data is sometimes exposed through the \ac{gbfs}~\cite{GBFS}. In broad terms, \ac{gbfs} defines common feeds for sharing both fixed information and time-varying availability so that third parties can use the data for mapping, journey planning, and operational monitoring~\cite{GBFS}. \ac{gbfs} is intended as a ``live'' interface rather than a historical archive and explicitly cautions against uses that could enable individual tracking or privacy risks~\cite{GBFS}.

Complementarily, some studies use data generated indirectly by the system or by its interaction with the surrounding environment. For example, social-network data can be leveraged to quantify system popularity~\cite{Das2015} or to identify perceived enablers and barriers to \ac{bss} use~\cite{Serna2019}. More generally, increasingly available sensing technologies--including cameras and active sensors such as LiDAR--can detect and track cyclists, micromobility vehicles, pedestrians, and motor vehicles, generating trajectory and interaction data that can be used to infer behavioral patterns and compute surrogate indicators of conflict risk at specific locations~\cite{Sayed2013, Saunier2007}. In addition, vision-based approaches and dedicated bicycle counters support monitoring of cycling flows, exposure estimation, before-after evaluation of interventions, and forecasting of cycling demand~\cite{Rogers2000, Biondi2022, Holmgren2017}.

Beyond system-generated data, \acp{bss} research uses a wide range of exogenous data. Contextual data such as weather conditions, public events, and seasonality strongly shape travel behavior~\cite{Bean2021, Julio2022, Wang2021, Park2022}, while built-environment and socioeconomic layers, including land-use patterns, population density, employment distribution, and proximity to \ac{pt}, provide essential information on the spatial determinants of demand~\cite{FaghihImani2014, FaghihImani2016, MateoBabiano2016, FaghihImani2017, Sun2018, FaghihImani2017Empirical}. In addition, infrastructure planning, accessibility analyses, and routing algorithms routinely draw on street-network data, typically obtained from sources such as \ac{osm}~\cite{Li2020, Zochowska2021} or released by public administrations~\cite{GarciaPalomares2012}. The integration of these rich, diverse datasets is now a standard foundation for advanced, data-driven research and operational strategies in the field of \acp{bss}. 



\section[Main fields]{Main fields}

With this data landscape in place, the main applications of computational methods in \acp{bss} can be grouped into several established fields. These areas build on the descriptive and exploratory analyses commonly used to characterize spatiotemporal patterns, flows, and other system-level behaviors. However, that foundational exploratory work is not revisited here, as it is already discussed in Chapters~\ref{ch:bss_users_and_their_mobility_patterns} and \ref{ch:impacts_of_bss}. The following Subsections therefore focus on the operational domains in which computational methods play a direct role.

    \subsection{Demand forecasting}

    Among the diverse applications of computational methods in \acp{bss}, demand forecasting stands out as one of the most established and intensively studied. The ability to accurately predict future rentals and returns at individual stations is foundational not only for many day-to-day operational tasks—such as rebalancing, pricing, battery charging, and maintenance scheduling—but also for higher-level strategic planning, including the siting of new stations based on anticipated use~\cite{Liang2023}. However, the inherently volatile nature of micromobility demand, shaped by a complex interaction of temporal patterns, weather events, spatial dynamics, and unpredictable social phenomena, poses considerable challenges for prediction.

    Initial research in this field focused on classical statistical time-series models, notably \ac{arima} and \ac{sarima}~\cite{Yoon2012, Kaltenbrunner2010}. These approaches remain valuable for modeling regular, periodic behaviors such as daily commuting peaks, and they are valued for their interpretability. Yet, their core assumptions of linearity and stationarity restrict their ability to integrate many external variables or adapt to sudden, non-linear shocks, such as unanticipated weather disruptions. Recognizing these shortcomings, research gradually shifted to more flexible and powerful machine learning approaches.

    \ac{ml} regressors subsequently emerged as a dominant paradigm due to their ability to accommodate high-dimensional feature spaces and complex relationships. Random Forests~\cite{Wang2018, Ruffieux2017}, Gradient Boosting Machines~\cite{Li2015, Sathishkumar2020}, and Support Vector Regression~\cite{Sathishkumar2020} are frequently used in the literature and typically outperform classical statistical approaches. Apart from using the input time series, these models generally rely on a diverse set of additional inputs, including weather variables (temperature, precipitation, wind), temporal features (hour, weekday, season), and contextual indicators such as holidays or events.

    Despite their effectiveness, traditional \ac{ml} techniques have an important limitation: they are often structured to treat each station independently or aggregate data at a coarse level, which hinders their ability to model spatial dependencies and interactions among stations that are central in \acp{bss}. This recognition has raised increasing interest in \ac{dl} approaches better suited to spatiotemporal data. For instance, \acp{rnn}—especially \ac{lstm} networks and \acp{gru}~\cite{Wang2018, Zhang2018LSTM}—as well as \acp{cnn}~\cite{Ruffieux2017}, have been widely adopted to capture temporal and local spatial structures in demand. Hybrid models, which combine the strengths of \acp{cnn} for spatial feature extraction with \acp{rnn} for modeling temporal dependencies, further boost predictive performance~\cite{Mehdizadeh2022}. Moreover, some studies have explored the use of hybrid architectures with self-attention mechanisms to improve predictive accuracy~\cite{Ma2024}. In parallel, recent work has moved from purely point forecasts to probabilistic predictions that quantify uncertainty, since many operational decisions are risk-sensitive and benefit from forecast distributions or intervals rather than single expected values~\cite{Lim2024_HGA, Lim2022_DeepAR}.

    Building on these advances, the latest state-of-the-art approaches deploy graph-based \ac{dl} models to directly capture the networked nature of \acp{bss}. \acp{gnn}~\cite{Liang2023}, along with more specialized variants such as \acp{stgnn}~\cite{Feng2024}, attention-based \acp{gnn}~\cite{Polamarasetti2025}, and \acp{stgcn}~\cite{Xiao2021}, represent each station as a node and use network edges to encode spatial proximity or historical trip flows. This graphical representation allows information to be propagated across the network, enabling these models to learn complex interdependencies and collectively predict demand patterns more effectively than methods relying on independent or grid-based structures. 

    Finally, demand forecasting in \acp{bss} must also contend with two recurrent modeling challenges. First, many operational objectives depend on \emph{net demand} (departures minus arrivals), which motivates models tailored to discrete station counts. For example, \citet{Liu2021_ExcessDemand} treat departures and arrivals as separate count processes and then forecast their difference using a Skellam model; this also helps infer \emph{excess demand} that is not directly observed when a station is empty or full. Second, station-level series can be sparse and noisy in large systems. To improve robustness, some approaches aggregate stations into geographically and behaviorally coherent clusters and forecast demand at higher levels, which are typically more regular than individual-station series and can already be used for planning rebalancing routes~\cite{Li2015TrafficPrediction}.


    \subsection{Fleet rebalancing}

    Fleet rebalancing represents another critical operational challenge in \acp{bss}, closely linked to the demand forecasting applications discussed previously. Because users pick up and return bicycles in an uncoordinated fashion, stations can become either empty or full. This imbalance can prevent rentals and returns, degrading overall service quality and creating lost demand that is not directly observed in system logs. Patterns of depletion and surplus often emerge systematically within a city, with residential areas depleting in the morning and facing surpluses by the evening, and employment or leisure areas typically displaying the converse pattern~\cite{Chiariotti2018}. To address these dynamic imbalances, operators rely on rebalancing which has thus become a central focus for optimization and \ac{ai}-driven approaches.

    The initial research efforts in this area approached the challenge using classical operations research formulations, which laid a theoretical foundation that subsequent methods build upon. Specifically, the \ac{sbrp} models overnight relocations, targeting a desired inventory state for all stations before the system resumes active demand. The \ac{dbrp} extends this formulation to situations where user activity continues while rebalancing is underway~\cite{Raviv2013}. Both the \ac{sbrp} and \ac{dbrp} are commonly formalized as mixed-integer (linear) programs or as pickup-and-delivery vehicle routing problems, with decision variables encoding the sequence of sites visited by each rebalancing vehicle and the corresponding loading or unloading of bicycles at each stop~\cite{Raviv2013, Chemla2013, DellAmico2014}.

    To solve these formulations, exact optimization techniques--including branch-and-cut, column generation, and related decomposition strategies--have been widely applied. Such approaches permit optimal solutions for small and medium-sized problems and provide valuable benchmarks for strategic planning or methodological evaluation~\cite{Chemla2013, DellAmico2014}. However, as system scale and complexity grow, the number of possible route and loading combinations increases combinatorially, making exact solution methods computationally prohibitive for large-scale or real-time rebalancing scenarios~\cite{Schuijbroek2017}.

    Recognizing the computational limitations, subsequent research has shifted toward approximate solution methods, primarily through decomposition strategies and metaheuristics, which in some cases are employed in tandem. Decomposition strategies, such as the ``cluster-first, route-second'' heuristic~\cite{Schuijbroek2017}, break the global problem into smaller, independent sub-problems by partitioning the system into zones. This enhances computational tractability, yet can induce sub-optimality at the boundaries between clusters. Complementing these, metaheuristics such as \acp{ga}~\cite{Li2016_Multiple}, Artificial Bee Colony~\cite{Shui2018}, Variable Neighbourhood Search~\cite{Kloimullner2014, Raidl2013}, and Adaptive Large Neighbourhood Search~\cite{Liu2025} provide flexible frameworks that explore the solution space and incorporate mechanisms such as \acp{ga}'s mutations to escape local optima. Moreover, some studies have combined heuristics with predictive models~\cite{He2021}, using demand forecasts to prioritize stations likely to face imminent shortages or surpluses. 
    
    Related work also emphasizes that rebalancing decisions must be made under uncertainty, since realized station flows often deviate from forecasts. For example, a two-stage stochastic model first chooses a morning allocation plan and then, once actual demand unfolds, computes a second-stage rebalancing action to correct the initial plan~\cite{Cavagnini2018}. Complementarily, the dynamic approach of Chiariotti et al.~\cite{Chiariotti2018} represents departures and arrivals as stochastic events, using the resulting probabilities to estimate the risk that a station will become empty or full and to trigger rebalancing proactively.

    Building on predictive and simulation-based advances, the literature has more recently embraced \ac{rl} as a promising paradigm for rebalancing. \ac{rl}-based formulations treat the \ac{bss} as a sequential decision-making environment, with rebalancing vehicles acting as agents that learn policies mapping system states—such as station inventories, predicted future demand, and vehicle locations—to relocation actions~\cite{Yin2024}. Deep \ac{rl} algorithms, including Deep Q-Networks and actor--critic architectures, have demonstrated an ability to outperform static heuristics in realistic, simulated environments~\cite{Yin2024, Liang2025}. This blend of learning and simulation enables adaptive policies that can respond dynamically to fluctuating network conditions, linking model prediction, decision-making, and feedback in an integrated loop.

    Complementing operator-based strategies, user-based rebalancing has emerged as an innovative approach that leverages the collective power of system users to improve inventory distribution. Instead of relying exclusively on vehicle fleets and staff, operators can offer incentives--credits, discounts, or gamified rewards--to prompt users to pick up bikes from stations that are oversupplied or return them to stations at imminent risk of depletion~\cite{Fukushige2022, Chiariotti2020}. Data-driven models aid in identifying the times and locations where such incentives are most likely to be effective~\cite{Chiariotti2020}, while optimization and \ac{rl}-based strategies can dynamically set incentive levels under operational and budget constraints~\cite{Cheng2021, Pan2019}.
    
    Although user-based rebalancing cannot fully substitute for operator-mediated relocations--particularly during severe \ac{bss} usage peaks--it provides a scalable and cost-effective way to supplement traditional rebalancing, reduce operational costs and emissions, and crowdsource part of the redistribution effort.   

    
    \subsection{Infrastructure planning}

    When discussing infrastructure in the context of \acp{bss}, attention often turns to decisions about station planning, including the number, location, type, and capacity of stations. Because this thesis devotes Part~\ref{pt:data-driven-station-planning} to this topic, with a comprehensive literature review provided in Chapter~\ref{ch:station_planning_in_barcelona}, station planning is not examined in detail here. Instead, this subsection focuses on the city’s cycling infrastructure, which plays a crucial role in enabling and encouraging \ac{bss} use~\cite{Moran2025, Otero2018}. 
    
    \subsubsection{Bike lanes}

    Bike lane planning models aim to design connected, accessible, and cost-effective networks that facilitate safe and convenient cycling. While high-level planning guidelines emphasize safety, connectivity, and comfort~\cite{Yanocha2018}, a growing body of modeling research supports these goals with formal methods.

    % Qualitative approaches
    Early approaches to bike lane planning relied largely on qualitative assessments or heuristic criteria, often grounded in expert judgement, safety audits, or stated- and revealed-preference surveys~\cite{Dondi2011, Tilahun2007, Sener2009}. These methods offered valuable insight into cyclists' perceptions but are often work-intensive and, in some cases, suffered from limitations such as working with small or context-specific samples~\cite{Tilahun2007}. As a response to these shortcomings, analytical optimization models emerged. While such models offered a more systematic approach, they traditionally relied on simplified assumptions, as in some cases they used only \ac{od} demand data without embedding the actual road network, ignoring route-choice behavior or network continuity constraints~\cite{Lin2011}. Additionally, the computational complexity of exact, mixed-integer formulations limited their scalability. 

    % GIS and MCDM approaches
    The increasing availability of digital spatial data and open-source routing tools introduced a new generation of data-driven approaches. \ac{gis}-based \ac{mcdm} frameworks have been used for evaluating potential interventions, combining layers such as population density and \ac{pt} proximity to identify high-priority corridors~\cite{Rybarczyk2010, Guler2021a}. These methods offered transparency and flexibility but still lacked the capacity to capture network effects, emergent route choices, or systemic continuity constraints.

    % Route-choice modelling: Stated-preference
    A key limitation of approaches based only on \ac{od} demand or static \ac{gis} layers is that they do not observe the street-level routes cyclists actually follow, and therefore cannot directly account for route-choice behavior when prioritizing links and corridors. One way to address this is to estimate cyclists' preferences over route attributes using stated-preference route-choice surveys, in which respondents choose between hypothetical route alternatives. These surveys quantify trade-offs between route attributes such as distance, surface quality, and cycling infrastructure, and can inform infrastructure design~\cite{Reckermann2024, Stinson2003}. Once obtained, these preference estimates can be operationalized in routing models that represent the relevant network attributes and minimize a generalized travel cost under different route types--fastest, most comfortable, quietest, and flattest--, yielding plausible routes between \ac{od} pairs even when detailed trajectories are unavailable~\cite{Hrncir2014}.

    % Route-choice modelling: Revealed-preference
    A major methodological leap occurred with the emergence of fine-grained mobility data, particularly \ac{gps} trajectories. These data allow planners to move beyond \ac{od} pairs and stated-preference route-choice surveys and observe the actual street-level routes cyclists follow. This shift has enabled a new class of data-driven bike-lane planning models in which infrastructure decisions are directly informed by observed cyclist movements~\cite{Liu2022, Bao2017}. These trajectory-based frameworks capture how cyclists respond to network gaps, detours, and continuity, and they link infrastructure design to predicted route choices rather than relying on heuristic criteria or simplified demand assumptions. 
    Notably, Yang and Mesbah~\cite{Yang2013} discuss the complementary roles of stated- and revealed-preference data for route-choice modelling, and highlight that while revealed-preference data better reflects realised choices, it can be noisy when obtained through route recall rather than direct tracking.
    
    % Network-science research
    In parallel, network-science research has examined how bicycle networks grow and what makes them effective. Using a synthetic network-growth model applied to dozens of cities, studies such as the one by Szell et al.~\cite{Szell2022} show that new bike lanes initially provide limited benefit, but connectivity and directness improve rapidly once the network surpasses a critical size. Their approach simulates different strategies for expanding a bicycle network and finds that coherent, well-connected growth consistently outperforms scattered additions. At a finer level, Vybornova~\cite{Vybornova2021} proposes an automated method based only on network topology to detect and rank discontinuities such as missing links and intersection breaks. Other network-science research examines how bicycle networks can be expanded most effectively using actual cycling mobility data. Olmos et al.~\cite{Olmos2020} combine \ac{od} matrices and \ac{gps} mobility data with network analysis to estimate where cycling demand is likely to be highest, then apply a percolation-based connectivity rule to select a small set of links that, once built, would form a city-wide connected backbone. In this approach, practical ``next investments'' are obtained by removing segments that already exist or are planned, leaving the missing links that would most improve overall connectivity. 
    

    \subsubsection{Shared-use paths}

    Limited public space and rising demand for active travel have increased interest in shared-use paths, which are off-road corridors that often have painted lanes, where pedestrians and cyclists share space. For these facilities, the key question is therefore not only ``how many users fit,'' but ``how often do they get in each other's way.'' Because shared-use-path performance is shaped mainly by user interactions, computational work often translates meetings, passings, and delayed passings--when overtaking is impeded--into measurable metrics for planning and design.

    A first, widely used line of work develops \ac{los} or Hindrance models that predict interaction rates from observable conditions--most notably pedestrian and bicycle volumes, path width, and markings--and aggregate them into a single measure of operational quality. This way of modeling traces back to early \ac{los} formulations for pedestrian--bicycle facilities~\cite{Botma1995} and later work explicitly grounded in user perceptions~\cite{Hummer2005}. Building on this logic, Nikiforiadis et al.~\cite{Nikiforiadis2020methodology} propose an event-based framework based on field measurements and user surveys to predict the frequency of different user interactions. By employing regression trees and applying \ac{ahp} to assign weights to each event type, their model produces a \ac{los} measure that more accurately reflects user-perceived discomfort, highlighting delayed passings as particularly disruptive.

    Beyond \ac{los} scores, recent work uses fixed cameras with automated tracking to extract pedestrian and cyclist trajectories and compute surrogate safety measures, enabling safety assessment without relying on relatively rare crash records. For example, Beitel et al.~\cite{Beitel2018} present an automated video-based conflict-analysis workflow for cyclist-pedestrian interactions. Finally, user-experience research shows why computational indicators should be interpreted alongside perception: Delaney et al.~\cite{Delaney2017} argue that focusing only on visible and observable conflict events can miss subtler aspects of how shared paths feel to users, and they report that pedestrians and cyclists may prefer different measures (e.g., segregation versus non-segregation), motivating clear guidance such as simple codes of conduct and ``keep left'' rules.

    
    \subsubsection{Safety-oriented interventions and evaluation}

    % Slow-speed zones
    Safety-driven cycling and micromobility infrastructure planning increasingly combines (i) targeted interventions that reduce conflict risk and (ii) computational and sensing-based methods to prioritize, design, and evaluate them. A first class of measures focuses on managing speeds in environments with high exposure to vulnerable road users. Low-speed zones are typically prioritized using crash histories and exposure measures, often complemented with contextual and demographic indicators to target locations where speed reduction is expected to yield the largest safety benefits~\cite{GRSF_LSZGuide}. In shared micromobility, these policies can be operationalized digitally: \ac{gps}-based geofences allow operators to cap device speeds or disable riding in sensitive areas~\cite{Hall2024, ChargeAndGo2025}. However, evidence on safety impacts remains mixed; for example, a case study at Texas University found no statistically significant reduction in e-scooter crash rates or injuries associated with campus slow-zone geofences~\cite{Hall2024}.

    A second major focus is intersection safety, since a large share of severe bicycle crashes occur at intersections~\cite{KeitaShindgikar2024, Alluri2017}. Signal-control countermeasures aim to reduce turning conflicts and improve cyclist visibility. One widely used treatment is the \ac{lbi}, which gives cyclists a short head start at signalized intersections by turning the bicycle indication green about 3--7~seconds before parallel vehicle movements~\cite{KeitaShindgikar2024}. This early release helps cyclists establish position in the intersection and mitigates conflicts such as ``right-hook'' crash events.

    Computational methods support \emph{ex ante} evaluation of these treatments by simulating potential timing plans and comparing outcomes across scenarios. Microsimulation has been used to test alternative bicycle signal operations (e.g., \acp{lbi}, exclusive bicycle phases, or bike scrambles) and quantify operational trade-offs such as delay, throughput, and \ac{los} before even implementing them~\cite{Kothuri2018}. Moreover, safety-oriented studies can combine simulation with surrogate conflict measures to approximate how timing changes affect vehicle--bicycle conflict risk; for example, one study found fewer simulated conflicts under a 5-second \ac{lbi} when compared to standard signal timing, particularly at higher bicycle volumes~\cite{Oh2018}.

    On the other hand, the growing deployment of sensing technologies strengthens safety-oriented planning by providing high-resolution measurements of exposure and interactions where crash data are sparse. Camera-based computer vision can track cyclists and motor vehicles at intersections, detect and characterize vehicle-bicycle conflicts, and compute surrogate safety measures to diagnose risky movements and evaluate treatments such as signal timing changes or protected designs~\cite{Sayed2013, Saunier2007}. Complementarily, video-based bicycle counting provides continuous volumes that support exposure estimation and before and after evaluation of safety interventions~\cite{Rogers2000}. Beyond fixed sensors, bicycle-mounted sensing (e.g., low-cost LiDAR) has also been explored to detect nearby vehicles and support early accident warning aimed at improving cyclist safety~\cite{Blankenau2018}.

    Finally, thermal safety is an emerging concern in a context of global warming, since extreme heat can deter riding and elevate health risks. Recent computational work therefore combines bike-sharing usage data with high-resolution thermal data to quantify network-wide heat exposure and prioritize corridor-level cooling where it delivers the largest reductions. Cabrera et al.~\cite{Cabrera2025}, for example, overlay 4.76~million New York City \ac{bss} trips onto 10~m \ac{utci} fields generated with the a urban microclimate model and show that exposure is highly concentrated on a small fraction of street segments, implying that targeted ``cool corridor'' interventions can yield disproportionate benefits. Complementarily, Zeng et al.~\cite{Zeng2025} construct diurnal exposure and \emph{excess}-exposure indicators from shared-bike trajectories and diurnal land-surface-temperature dynamics to identify \emph{when} and \emph{where} riders accumulate the highest thermal burden, supporting time-sensitive measures. To support implementation at scale, image-based assessments can directly measure shade along cycling links: Cao et al.~\cite{Cao2024} derive a street-level Shade Index for bike lanes from street-view imagery using \ac{dl}, enabling network-wide assessment of shading, while Chen et al.~\cite{Chen2025StreetView} similarly use street-view indicators to examine how streetscape conditions influence cycling behavior during extreme heat.
